\subsection{A genuine witness with the stated signed finite normalization}

\begin{result}{E.196}{Proved}
{The parity-one witness can retain zero full shear and the original two finite frames}
There is a genuine full positive gentle element \(g\) with
\[
       \lambda(g)=1,\qquad \rho(g)=0,\qquad
       T_g=5I,\quad W_g=I\pmod{12}.
\tag{E.196.1}
\]
It fails the full pentagon and the common CM covariance
of E.190.  Its restrictions to the actual relative
pronilpotent congruence quotients are the explicitly
transformed restrictions described below.
This proves the existence of a full witness with the
historical displayed signed normalization.
It does not identify that element with an unspecified
choice in the historical A.2 relative-lifting family,
or retain an involution condition after the final
arithmetic correction.
\end{result}

\begin{proof}
Use the single arithmetic \(\sigma\) of E.96 and the
single full positive gentle \(h\) constructed in E.97
from that same \(\sigma\).  Thus \(\rho(\sigma)=0\),
\(\sigma(r_*)=-r_*\), and, with the literal frames,
\[
        T_\sigma=5I,\qquad
        W_\sigma\bmod3\in\{S,S^{-1}\},\qquad
        W_\sigma=I\pmod4 .
\tag{E.196.2}
\]
The finite action of \(h\) on every \(Q_N\) is
identity, its \(r\)-frame has image one, and its
\(s\)-frame has the image of \(\beta_\sigma\).
In particular its projective action at three is
identity.  E.95.6 identifies that action with
conjugation by \(V_h=U^{8\rho(h)}\).
The center of \(\operatorname{PSL}_2(\mathbb F_3)=A_4\)
is trivial; a power of this unipotent is scalar only
when it is identity.  Hence \(\rho(h)=0\pmod3\),
and \(V_h=I\pmod{12}\).
The exact stated frame images therefore give
\[
                       T_h=I,\qquad W_h=W_\sigma .
\tag{E.196.3}
\]

The class-two proof in E.95 shows that
\(\rho(h)\in2\widehat{\mathbb Z}\) for every
positive gentle element; it uses the two-cycle
and hexagon, without the pentagon.
The pure-frame law likewise makes \(\rho\)
additive on that group.
Set \(k=h\sigma^{-1}\), and put
\[
                   d=\rho(k)=\rho(h)
                        \in6\widehat{\mathbb Z}.
\tag{E.196.4}
\]
The homomorphism laws for the signed \(T,W\)
in E.95 yield
\[
                         T_k=5I,\qquad W_k=I .
\tag{E.196.5}
\]
The remaining correction must remove \(d\)
without losing either signed matrix.

By the exact Kummer image in E.181 choose
\(\tau_0\in G_{\mathbb Q^{\rm ab}}\) with
\(\kappa_2(\tau_0)=d\).
Because \(r_*^2=\sqrt3/2\) belongs to
\(\mathbb Q^{\rm ab}\), its image on \(r_*\)
is a sign.  Multiplication by \(\sigma\),
if necessary, makes that sign positive
without changing \(\kappa_2\).  Denote the
result by \(\tau_1\).

We check the full finite frames, not just
their bits.  The three-torsion of
\(E_3:Y^2=X^3+1\) is defined over
\[
                 \mathbb Q^{\rm ab}(2^{1/3}).
\tag{E.196.6}
\]
Indeed its third-division polynomial is
\(3X(X^3+4)\).  At \(X=0\) the ordinates
are \(\pm1\); at the other roots they are
\(\pm\sqrt{-3}\).  Roots of unity and those
ordinates lie in \(\mathbb Q^{\rm ab}\),
and the remaining abscissae are generated
by \(2^{1/3}\).
Since \(d=0\pmod3\), \(\tau_1\) fixes
all these points.  It also fixes the
equianharmonic quadratic twist because
it fixes \(r_*\), by E.95.4.
The equianharmonic signed frame is therefore
identity modulo three, and identity modulo
four by the same E.95 calculation:
\[
                             T_{\tau_1}=I .
\tag{E.196.7}
\]

On the harmonic fibre, fixing \(r_*\) fixes
the four third-division abscissae, but this
alone gives only \(W_{\tau_1}=\pm I\pmod3\).
The modulo-four matrix is again identity.
This residual sign must be removed separately.
By E.196.2,
\[
              T_{\sigma^2}=I,\qquad
              W_{\sigma^2}=-I\pmod3,\quad
              W_{\sigma^2}=I\pmod4,
              \qquad\kappa_2(\sigma^2)=0 .
\tag{E.196.8}
\]
Multiply \(\tau_1\) by \(\sigma^2\) exactly
when its harmonic sign is negative.  The
resulting arithmetic \(\tau\) satisfies
\[
             \kappa_2(\tau)=d,\qquad
                       T_\tau=W_\tau=I\pmod{12}.
\tag{E.196.9}
\]
The character and frame laws now show that
\(g=k\tau^{-1}=h\sigma^{-1}\tau^{-1}\)
has precisely E.196.1.
E.95 gives \(\Delta(g)=1\); E.191 therefore
shows that \(g\) cannot satisfy the pentagon.
More directly, its two signed bits differ,
so E.191.20 cannot hold for one exponent
\(j\in\mathbb Z/6\).  Thus the failure is
already at the stated common covariance
conclusion.

For clarity, record what the arithmetic
correction does to the relative data.
For every \(12\mid N\), let
\(\Gamma_{N,p}\) be E.97's actual relative
pro-\(p\) factor above \(Q_N\).
All the actions below preserve these
actual congruence kernels, by the finite
congruence covariance used in E.95--E.97.
Writing compositions in their action order,
\[
 \begin{split}
  g|_{\Gamma_{N,2}}&=\tau_2^{-1},\\
  g|_{\Gamma_{N,p}}&=\sigma_p^{-1}\tau_p^{-1}
                              \qquad(p\text{ odd}).
 \end{split}
\tag{E.196.10}
\]
Their common finite action is that of
\(\tau^{-1}\) on \(Q_N\), so they form one
actual compatible fibre-product action.
The formula is simply the composition
of E.97's hybrid \((\sigma_2,1_{\rm odd})\)
with the two genuine arithmetic actions.
Each accessory frame is the image of the
same full frame of \(g\), as follows from
the literal pure-frame composition law.
Thus no new lifting or compactness assertion
is concealed in E.196.10.  These are
transformed relative data; involutivity of
their negative companions and identity of
the old relative action are not asserted.
\end{proof}
